\bookchapter{제3장: 스마트폰에 “생명선” 연결하기 - 네트워크}

\chapterintro{인터넷이 없으면 스마트폰은 절반만 “똑똑”합니다. 네트워크는 현대 스마트폰의 생명선이며, 이를 이해하고, 잘 활용하고, 보호하는 것은 모든 스마트폰 사용자의 필수 과정입니다.}

\booksection{인터넷 연결에 혼란 없기 - Wi-Fi와 데이터의 “길”과 “차” 구분하기}

많은 사람들의 스마트폰 화면 상단에는 Wi-Fi와 신호 아이콘이 동시에 표시되지만, 둘의 본질적 차이에 대해서는 혼란스러워합니다. 가장 고전적인 오해는:\textbf{Wi-Fi에 연결되어 있는데 페이지를 열 수 없음}- 이는 라우터가 인터넷에 연결되지 않았기 때문이며, 스마트폰이 자동으로 데이터로 전환되지 않았기 때문입니다. 이 문제를 완전히 이해하려면 “길과 차”의 인지 모델을 구축해야 합니다.

\begin{tcbinfo}{핵심 비유: 길과 차}
인터넷을 도시라고 상상해 보세요.\textbf{“길”}은 스마트폰에서 인터넷까지의 통로이고,\textbf{“차”}는 그 위를 달리는 데이터입니다. Wi-Fi와 데이터는 서로 다른 두 길이지만, 데이터(차) 자체는 동일합니다.

\end{tcbinfo}

\booksubsection{본질적 차이: 로컬 통로 vs 원격 고속}

Wi-Fi는 본질적으로\textbf{근거리 네트워크 기술}로, 집 아래에 무선으로 깔린 길과 같아서 라우터에서 스마트폰으로 데이터를 전송하는 역할을 합니다. 하지만 broadband 요금이 미납되었거나 광케이블이 끊어지면, 길은 그대로 있지만 외부의 “고속도로”가 끊어져 어디에도 갈 수 없습니다. 반면 모바일 데이터는 이동통신 사업자의 4G/5G 기지국을 통해 전송되어, 사업자의 고속 전용 라인에 직접 올라타는 것과 같아서 신호만 있으면 인터넷에 직접 연결할 수 있습니다.

\booksubsection{요금과 시나리오: 고정 월정액 vs 사용량 기반 요금}

Wi-Fi는 일반적으로 일회성 투자(broadband 월 요금)이며, 전기만 공급되면 가족 모두가 무제한 사용할 수 있어 영화 다운로드, 게임 등 대용량 시나리오에 매우 적합합니다. 반면 데이터는 사용량에 따라 요금이 부과되며, 요금제를 초과하면 속도가 느려지거나 추가 요금이 발생할 수 있어 지하철, 출장 등 이동 시나리오에서 간단한 작업을 처리하는 데 더 적합합니다.

\begin{tcbpractice}{시도해 보기: 데이터 요금제 확인}
스마트폰의 [설정] > [모바일 네트워크] > [데이터 관리]를 열고, 이번 달에 남은 데이터 양과 어떤 앱이 데이터를 가장 많이 사용하는지 확인합니다. 상황을 파악하고 있어야 외출해도 당황하지 않습니다.

\end{tcbpractice}

\booksubsection{보안 고려사항: 공용 Wi-Fi는 숨겨진 “피싱장”}

보안 측면에서 모바일 데이터는 일반적으로 공용 Wi-Fi보다 안전합니다. 모바일 데이터는 이동통신 사업자가 암호화하지만, 공용 Wi-Fi는 해커가 “가짜 핫스팟”을 구축하여 암호화되지 않은 데이터를 가로채기 쉽습니다. 그러나 HTTPS가 전면적으로 보급된 현재, 공격자가 가로챌 수 있는 것은\textbf{방문한 웹사이트(도메인)}와\textbf{암호화되지 않은 HTTP 콘텐츠}입니다. 위챗 메시지, 은행 앱의 로그인 비밀번호 등 HTTPS로 암호화된 데이터는 공격자가\textbf{직접 읽을 수 없습니다}- 더 복잡한 SSL 스트리핑 공격을 수행하고 사용자가 인증서 경고를 무시하는 경우는 예외입니다. 그럼에도 불구하고 외출 중 결제 작업을 할 때는 모바일 데이터를 직접 사용하는 것이 가장 안전한 선택입니다.

\textbf{보안리마인더}\par

공항, 쇼핑몰, 카페 등 공공장소에서는 의심스러운 이름의 무료 Wi-Fi(예: “Free\_WiFi\_2024”)에 연결하지 마세요. 연결했더라도 은행 앱에 로그인하거나 비밀번호를 입력해서는 절대 안 됩니다.

\booksubsection{스마트폰에서 판단: 상태 표시줄의 “암호” 읽기}

스마트폰이 현재 어떤 방식으로 인터넷에 연결되어 있는지 어떻게 알 수 있나요? 화웨이 스마트폰을 예로 들면, 상태 표시줄에서 이동통신 사업자 신호 옆 또는 WLAN 신호 옆에\textbf{양방향 화살표 아이콘}이 나타나면 해당 방식으로 인터넷에 연결 중임을 의미합니다. 둘이 동시에 켜져 있으면 네트워크 가속이 작동 중이며, 스마트폰이 두 통로를 동시에 사용하고 있다는 뜻입니다.

iPhone 사용자에게는 더 직관적입니다. Wi-Fi 아이콘이 켜져 있으면 Wi-Fi를 사용 중이고, 꺼져 있을 때 우측 상단의 신호 막대를 보면 데이터를 사용 중임을 알 수 있습니다. 일부 안드로이드 스마트폰은 상태 표시줄에 “5G” 또는 “4G” 문자를 표시하여 현재 모바일 네트워크 유형을 직접 알려줍니다.

\booksection{모바일 네트워크(4G/5G)와 데이터}

모바일 네트워크는 언제 어디서나 인터넷에 연결할 수 있는 기반입니다. 2G에서 5G로, 속도는 수천 배 빨라졌지만 원리는 사실 비슷합니다.

\booksubsection{네트워크 규격:}

\begin{itemize}
  \item \textbf{4G LTE}: 현재 가장 널리 보급되었으며, 속도가 충분하고(이론상 최대 100Mbps+), 전력 소비가 상대적으로 낮습니다.
  \item \textbf{5G}: 속도가 빠르지만(이론상 최대 10Gbps), 커버리지가 아직 완전하지 않고 전력 소비가 높습니다. 대용량 파일 다운로드, HD 영상 시청에 적합합니다.
  \item \textbf{5G SA vs NSA}: SA는 독립 구성망, NSA는 4G 기지국에 의존합니다. SA 경험이 더 좋지만, 이동통신 사업자와 스마트폰 모두에서 지원해야 합니다.
\end{itemize}

\booksubsection{데이터 관리:}

\begin{itemize}
  \item \textbf{데이터란}: 스마트폰으로 인터넷을 사용할 때 발생하는 데이터 양입니다. 영상 시청, 이미지 구경이 데이터를 가장 많이 소비하며, 텍스트와 음성은 상대적으로 적게 듭니다.
  \item \textbf{데이터 절약 팁}: 설정에서 "데이터 절약 모드"를 활성화하고, Wi-Fi 환경에서 앱 업데이트나 HD 영상 시청을 합니다.
  \item \textbf{신호 막대 $\neq$ 속도}: 신호가 강하다는 것은 연결이 좋다는 의미일 뿐, 속도는 기지국이 얼마나 바쁜지, 기지국에서 얼마나 떨어져 있는지에 따라 다릅니다.
\end{itemize}

\textbf{중요 경고: 출국 전 [데이터 로밍]을 꼭 끄세요!}\par

해외에서 모바일 데이터를 사용하면 엄청난 요금이 발생할 수 있습니다! 출국 전 반드시 설정에서 데이터 로밍을 끄거나, 현지 데이터 SIM 카드/국제 로밍 요금제를 미리 구매하세요.

\booksection{Wi-Fi: 연결, 공유 및 보안}

Wi-Fi는 가정과 사무실에서 인터넷을 사용하는 주력 수단으로, 속도가 빠르고 데이터를 소모하지 않습니다. 하지만 Wi-Fi도 계속 발전하고 있어, 표준 간 차이가 큽니다.

\booksubsection{연결Wi-Fi:}

[설정] > [WLAN]에서 네트워크를 선택하고 비밀번호를 입력합니다. 이제 많은 스마트폰이 QR 코드 스캔 연결을 지원하여 더 편리합니다.

\booksubsection{공유비밀번호:}

안드로이드와 iOS 모두 연결된 스마트폰이 비밀번호를 친구에게 빠르게 공유할 수 있습니다. 안드로이드는 QR 코드를 생성하고, iOS는 가까이 대면 자동으로 공유할 수 있습니다.

\booksubsection{Wi-Fi 표준의 진화:}

Wi-Fi 표준은 IEEE가 제정하며, 세대마다 큰 향상이 있습니다:

\begin{center}
\begin{tabularx}{\linewidth}{|X|X|X|X|X|}
\hline
\textbf{표준} & \textbf{프로토콜} & \textbf{주파수 대역} & \textbf{최고 속도} & \textbf{특징} \tabularnewline
\hline
Wi-Fi 4 & 802.11n & 2.4/5GHz & 600Mbps & 오래된 표준, 커버리기는 괜찮음 \tabularnewline
Wi-Fi 5 & 802.11ac & 5GHz & 3.5Gbps & 현재 가장 보급되어 있고, 속도도 충분함 \tabularnewline
Wi-Fi 6 & 802.11ax & 2.4/5GHz & 9.6Gbps & 다중 기기 성능이 강하고 지연이 낮음 \tabularnewline
Wi-Fi 6E & 802.11ax & 2.4/5/6GHz & 9.6Gbps & 6GHz 주파수 대역 추가, 혼잡 없음 \tabularnewline
Wi-Fi 7 & 802.11be & 2.4/5/6GHz & 46Gbps & 초고속 전송, 다중 링크 집계 \tabularnewline
\hline
\end{tabularx}
\end{center}

\begin{itemize}
  \item \textbf{Wi-Fi 6의 핵심 기술}: OFDMA(직교 주파수 분할 다중 접속)는 여러 기기가 동시에 전송할 수 있게 하고, MU-MIMO(다중 사용자 다중 입출력)는 다중 기기 동시 효율을 높이며, TWT(타깃 웨이크 타임)는 기기를 더 전력 절약되게 합니다.
  \item \textbf{Wi-Fi 6E}: Wi-Fi 6에 6GHz 주파수 대역을 추가하여 더 많은 채널을 제공하고 간섭을 줄입니다. 스마트폰과 라우터 모두에서 지원해야 합니다.
  \item \textbf{Wi-Fi 7}: MLO(다중 링크 운영)를 도입하여 여러 주파수 대역을 동시에 사용하여 데이터를 전송하고 지연을 크게 줄입니다. 320MHz 초광대역 채널로 속도가 도약합니다. 스냅드래곤 8 Elite Gen5, 미디어텍 9400 등 플래그십 칩 플랫폼이 모두 Wi-Fi 7을 지원하며, Xiaomi 17, 화웨이 Mate 80, 삼성 S26 등 모델은 Wi-Fi 7 네트워크를 사용할 수 있습니다(Wi-Fi 7 라우터와 함께 사용해야 함).
  \item \textbf{2.4GHz vs 5GHz}: 2.4GHz는 벽 투과가 좋지만 속도가 느리고 간섭이 많습니다. 5GHz는 속도가 빠르지만 벽 투과가 약합니다. Wi-Fi 6E/7의 6GHz 주파수 대역은 속도가 가장 빠르지만 벽 투과가 가장 약해, 같은 방 안에서 사용하기에 적합합니다.
\end{itemize}

\booksubsection{공용 Wi-Fi 위험:}

공용 Wi-Fi는 해커가 데이터를 도청하거나 피싱 핫스팟을 설치할 수 있습니다. 수칙:

\begin{itemize}
  \item 은행 계좌 로그인, 비밀번호 입력을 피하세요.
  \item VPN으로 연결을 암호화하는 것이 가장 좋습니다.
  \item 낯선 Wi-Fi에 자동으로 연결하지 마세요.
  \item HTTPS 웹사이트(URL 입력란에 자물쇠 아이콘이 있는)를 사용하세요.
\end{itemize}

\booksection{블루투스: 단거리 무선 통신의 만능 선수}

블루투스는 스마트폰에서 가장 자주 사용하는 무선 기술 중 하나일 것입니다 - 무선 이어폰, 스마트워치, 차량 시스템, 키보드 마우스 모두 블루투스로 연결됩니다.

\booksubsection{기술 원리:}

블루투스는 2.4GHz ISM 대역에서 작동하며, 주파수 도약 확산 기술(FHSS)을 사용하여 79개 채널에 데이터를 분산하여 빠르게 전환 전송함으로써 간섭을 효과적으로 회피합니다. 블루투스의 통신 거리는 세 등급으로 나뉩니다: Class 1 약 100미터, Class 2 약 10미터(스마트폰에서 주로 사용), Class 3 약 1미터.

\booksubsection{버전 진화:}

\begin{itemize}
  \item \textbf{블루투스 4.0}: 저전력 블루투스(BLE)를 도입하여, 손목 밴드, 심박계 같은 소형 배터리 기기도 무선으로 연결할 수 있게 되어 배터리 수명이 며칠에서 몇 달로 늘어났습니다.
  \item \textbf{블루투스 4.2}: 전송 속도와 프라이버시 보호가 향상되었으며, IPv6 직접 연결을 지원합니다.
  \item \textbf{블루투스 5.0}: 전송 속도 2배(2Mbps), 전송 거리 4배(이론상 300미터), 방송 능력 8배. 듀얼 오디오(이어폰 두 개 동시 연결) 지원.
  \item \textbf{블루투스 5.1}: 방향 탐지 기능(AoA/AoD) 도입, 위치 정확도가 센티미터 수준에 도달.
  \item \textbf{블루투스 5.2}: LE Audio 신표준, LC3 코덱 지원(음질이 더 좋고 전력 소비가 더 낮음), 다중 스트림 오디오, 방송 오디오.
  \item \textbf{블루투스 5.3/5.4}: 전력 소비와 연결 안정성을 더욱 최적화하고, 주기적 방송을 강화.
\end{itemize}

\textbf{LE Audio 보급 현황}: 2025-2026년, LE Audio는 빠른 보급 단계에 진입했습니다. 주요 브랜드의 신형 완전 무선 이어폰(예: 화웨이 FreeBuds Pro 5, 샤오미 Buds 6 Pro, 삼성 Galaxy Buds3 Pro 등)은 모두 LE Audio와 LC3 코덱을 지원하여 더 낮은 지연과 더 나은 음질을 제공합니다. 새 이어폰 구매 시, LE Audio 지원 제품을 우선 선택할 것을 권장합니다.

\booksubsection{활용 시나리오:}

\begin{itemize}
  \item \textbf{오디오 전송}: 무선 이어폰, 차량용 오디오. 블루투스 5.2의 LE Audio는 음질과 지연 모두 질적 도약을 가져왔습니다.
  \item \textbf{데이터 전송}: 근거리로 파일, 사진 전송. Wi-Fi Direct보다 속도는 느리지만 단순한 장점이 있습니다.
  \item \textbf{기기 페어링}: 스마트워치, 손목 밴드, 키보드 마우스, 게임 패드.
  \item \textbf{위치 서비스}: 블루투스 비콘(Beacon)은 실내 내비게이션, 물품 추적(예: Apple AirTag)에 사용됩니다.
  \item \textbf{스마트 홈}: 블루투스 메시 네트워크로 전구, 도어락 등을 제어.
\end{itemize}

\booksection{적외선 리모컨: 스마트폰이 만능 리모컨으로 변신}

적외선 리모컨은 가장 오래된 무선 제어 기술 중 하나이지만, 스마트폰에서는 여전히 유용합니다. 일부 안드로이드 스마트폰(예: 샤오미, 화웨이 일부 모델)은 적외선 발신기를 내장하여 TV, 에어컨, 셋톱박스 등 가전제품을 제어할 수 있습니다.

\booksubsection{작동 원리:}

적외선 리모컨은 적외선 발광 다이오드를 통해 특정 주파수(보통 38kHz)의 적외선 펄스 신호를 발신하고, 가전제품의 적외선 수신기가 이러한 펄스를 디코딩하여 해당 작업을 수행합니다. 브랜드마다 다른 인코딩 프로토콜을 사용하며, 일반적으로 NEC, RC5, RC6, SONY 등이 있습니다. 스마트폰의 적외선 앱에는 각 브랜드의 인코딩 라이브러리가 내장되어 있어, 브랜드와 모델을 선택하면 매칭됩니다.

\booksubsection{한계:}

\begin{itemize}
  \item \textbf{방향성}: 적외선은 직선으로 전파되므로, 기기를 향해 발사해야 하며 벽을 사이에 두고 사용할 수 없습니다.
  \item \textbf{짧은 거리}: 유효 거리는 일반적으로 5-8미터입니다.
  \item \textbf{애플 미지원}: iPhone에는 적외선 발신기가 없어 이 기능을 사용할 수 없습니다.
\end{itemize}

\booksection{NFC: 터치의 마법}

NFC(Near Field Communication, 근거리 무선 통신)는 매우 짧은 거리(10cm 이내)의 고주파 무선 통신 기술로, 작동 주파수는 13.56MHz이며 전송 속도는 최대 424Kbps입니다.

\booksubsection{통신 원리 및 모드:}

NFC는 RFID 기술에서 발전한 것으로, 세 가지 작동 모드를 지원합니다:

\begin{itemize}
  \item \textbf{읽기/쓰기 모드}: 스마트폰이 NFC 태그의 정보를 읽습니다. 예: NFC 스티커를 읽어 동작을 트리거.
  \item \textbf{P2P 모드}: 두 NFC 기기가 서로 데이터를 전송합니다. 예: Android Beam(중단됨).
  \item \textbf{카드 에뮬레이션 모드}: 스마트폰이 스마트카드를 에뮬레이트하며, 가장 일반적으로 사용되는 모드입니다 - 교통카드, 출입카드, 은행 카드가 모두 이 모드입니다.
\end{itemize}

\booksubsection{보안 메커니즘:}

\begin{itemize}
  \item \textbf{초단거리}: 10cm 내의 통신 거리 자체가 물리적 보안 보장으로, 원격 도청은 거의 불가능합니다.
  \item \textbf{보안 칩}: 스마트폰에는 독립적인 보안 칩(SE)이 내장되어 있거나 HCE(호스트 카드 에뮬레이션)를 사용하며, 결제 키는 보안 영역에 저장되어 스마트폰이 루팅되어도 추출할 수 없습니다.
  \item \textbf{지문/얼굴 인증}: 결제 시 생체 인식 확인이 필요하여, 스마트폰을 분실해도 도난 결제할 수 없습니다.
  \item \textbf{토큰화}: Apple Pay, 화웨이 Pay 등은 실제 카드 번호를 전송하지 않고 일회용 토큰을 사용하므로, 가로채도 재사용할 수 없습니다.
\end{itemize}

\booksubsection{대표적 활용 사례:}

\begin{itemize}
  \item \textbf{교통카드/출입카드}: 실물 카드 정보를 스마트폰에 저장하여 스마트폰으로 태그합니다.
  \item \textbf{모바일 결제}: Apple Pay, 화웨이 Pay 등은 은행 카드를 스마트폰에 연동하여 터치로 결제합니다.
  \item \textbf{기기 페어링}: NFC 터치로 블루투스 이어폰, 스피커를 빠르게 페어링.
  \item \textbf{본인 인증}: NFC 신분증 읽기, 전자 운전면허증 검증.
\end{itemize}

\booksection{개인 핫스팟: 스마트폰을 Wi-Fi 라우터로 변신}

자신의 데이터가 많을 때, 핫스팟을 켜서 태블릿, 컴퓨터에 공유할 수 있습니다.

\begin{itemize}
  \item \textbf{활성화 방법}: 설정 > 개인 핫스팟/네트워크 공유에서 활성화합니다.
  \item \textbf{주의사항}: 전력 및 데이터를 많이 소모하므로, 복잡한 비밀번호를 설정하고 연결 기기 수를 제한해야 합니다.
  \item \textbf{USB 테더링}: 데이터 케이블로 컴퓨터에 연결하며, 무선 핫스팟보다 더 안정적이고 빠릅니다.
\end{itemize}

\begin{tcbpractice}{직접 해보기: 핫스팟 설정 및 관리}
\begin{enumerate}
  \item [설정]을 열고 "개인 핫스팟" 또는 "네트워크 공유"를 찾아 활성화한 후, 이름과 비밀번호를 설정합니다.
  \item 다른 기기(예: 컴퓨터)로 이 Wi-Fi를 검색하여 연결합니다.
  \item 스마트폰에서 이미 연결된 기기를 관리하거나, 모르는 기기를 강제 종료할 수 있는지 확인합니다.
\end{enumerate}

\end{tcbpractice}
